\quellseite{683}
und einen Balken machen aus dem Splitter. Denn der Zorn ist ein junges Reis, der Haß aber ist ein Baum und ein großer Balken. Daher gehört auch dies unter das fünfte Gebot, wofür wir den Apostel Johannes als Gewährsmann haben [1 Joh. 3, 15.]: „Wer seinen Bruder hasset, der ist ein Todtschläger.“

Die siebente Todsünde ist die Trägheit, das ist die Verdrossenheit zum Guten, Faulheit, griechisch ἀργία. Diese ist zwiefach. Die eine ist in dem figürlichen (figurativo) und äußerlichen Gottesdienste, als da sind: die Kirche besuchen, das Wort hören, beten, lesen, betrachten, singen, und in solcher Weise ist dies Laster unter dem dritten Gebote verboten: „Du sollst den Feiertag heiligen“, wie genugsam offenbar ist. Die andere ist feiner und innerlicher und findet statt in dem ganzen Dienste GOttes und in jeder Art desselben, und in dieser Weise ist sie nichts Anderes als das Vertrauen auf die angefangene Gerechtigkeit und die Nachlässigkeit im Zunehmen; daß man stillsteht auf dem Wege GOttes, lauwarm ist, sicher wird, die Furcht GOttes fahren läßt. Von diesem Laster der Heuchler redet die ganze Schrift. Denn diese Trägheit macht sichere Werktheilige, vor deren Augen nicht

die Furcht GOttes ist, und „die Gerichte GOttes sind fern von ihnen“ [Ps. 10, 5.]. Mit diesen Leuten wird es hernach ärger denn vorhin [Luc. 11, 25. f.], denn der Teufel findet ihr Haus geschmückt, aber leer und sicher. Deshalb gehen die Teufel hinein und wohnen daselbst. O welch ein überaus erschreckliches Wort, daß die Teufel daselbst wohnen, weil die Menschen, welche derartig sind, nach ihrer angefangenen Gerechtigkeit viel verhärteter werden als zuvor, wie wir in der Erfahrung sehen, daß sie sich zu dem, was sie hören, nicht bewegen lassen, sondern glauben, das treffe sie nicht, und sind in Wahrheit ein Otterngezüchte geworden [Matth. 3, 7. 12.], eine Spreu, die bereitet ist für das ewige Feuer. Also findet sich diese Sünde bei allen Geboten, da man in einem jeglichen zunehmen soll. Aber ich weiß nicht, ob man sie beichten soll. Ich glaube, daß man es nicht thun solle, denn es ist ein geistliches Gebrechen, welches man allein GOtt eröffnen soll, der auch nur allein daselbst heilen kann.\footnote{In der Baseler, der Wittenberger und der Jenaer Ausgabe ist am Schlusse hinzugefügt: „Ende der Predigten über die zehn Gebote, welche der ehrwürdige Vater Martin Luther im Jahre 1516 und 1517 für das Volk in Wittenberg gehalten hat.“}

\chapter[D. Die zehn Gebote GOttes mit einer kurzen Auslegung]{D. Die zehn Gebote GOttes mit einer kurzen Auslegung ihrer Erfüllung und Uebertretung}

Von D. Martin Luther, Augustiner, gemacht.\footnote{Wie wir aus dem Anfang der unmittelbar vorhergehenden Predigt über die sieben Todsünden ersehen (Col. 1344), war durch die unendliche Vervielfältigung der Arten der Sünden aus der Beichte eine grausame Marter der Gewissen geworden. Deshalb schrieb Luther für die armen geplagten Seelen zu Anfang des Jahres 1518 eine Anweisung, wie man nach den heiligen zehn Geboten die Beichte kurz einrichten könne, beide in lateinischer und deutscher Sprache. Die lateinische Schrift hat den Titel: Instructio pro confessione peccatorum abbrevianda secundum decalogum. Autore D. M. Luth. Anno MDVIII. und findet sich in der lateinischen Wittenberger Ausgabe (1550), Tom. I, fol. 199; Jenaer (1579), Tom. I, fol. 175b; Erlanger, exeg. opp., Tom. XII, p. 219 und in der Weimarischen, Bd. I, S. 257. Die einzige lateinische Einzelausgabe dieser Schrift, welche in der Weimarischen Ausgabe angegeben ist, erschien erst im Mai 1520 bei Silvanus Ottmar in Augsburg unter einem Titel, der den der deutschen Ausgabe nachgebildet ist. Die deutsche Bearbeitung ist, wie sich durch Vergleichung der Schrift: „Kurze Form, die zehn Gebote, Glauben und Vater-Unser zu betrachten“ (Wald, St. Louiser Ausg., Bd. X, 153–159) ergibt, von Luther selbst; in letzterer Schrift hat er nur den ersten Theil ganz verändert, dagegen bei Anzeigung der Uebertretung und Erfüllung erweitert; im Uebrigen stimmt sie mit unserer Schrift von Wort zu Wort. „Eine Sonderausgabe der deutschen Bearbeitung von 1518 läßt sich nicht nachweisen, aber wir haben einen Abdruck innerhalb folgender Schrift Caspar Göttels: ‚Jhesus. Ein fast fruchtbar büchlein von Adams wercken, vnd gottes genade mit vnterricht wie recht beichten, busszen, vnd das hochwirdigst Sacrament selig tzu empfahen im Augustiner Closter tzu sandt Anne vor Eisleben dise heiligste fasten gepredigt vnd gegeben. 1518.‘ (Druck von Melchior Lotther in Leipzig.) Hier findet sich die Anleitung deutsch, Bl. Ciiija—Dijb, inmitten des fünften Capitels; sie hat also in den Fasten 1518 schon gedruckt vorgelegen. (Weim. Ausg.) In der Weimarischen Ausgabe werden unter dem von uns angegebenen Titel neue Einzelausgaben angeführt, deren keine in Wittenberg erschienen ist. Mehrere derselben sind ohne Zeitangabe, nämlich zwei bei Jobst Gutknecht in Nürnberg, eine bei Adam Petri in Basel, eine bei S. Otmar (Ottmar) in Augsburg, eine bei Jörg Nadler in Augsburg; im Jahre 1519 eine bei Wolfgang Stöckel in Leipzig, eine im Jahre 1520 bei S. Otmar zu Augsburg, und ebendaselbst 1522; eine bei Hans Froschauer zu Augsburg 1523. In den Sammlungen: in der Eislebenschen, Bd. I, Bl. 21; in der Altenburger, Bd. VIII, S. 894; in der Leipziger, Bd. IV, S. 86; in der Erlanger, Bd. 36, S. 145 und in der Weimarischen, Bd. I, S. 247. In den verschiedenen Ausgaben ist eine große Verschiedenheit in der Anordnung des Textes, welche daher rührt, daß in dem verloren gegangenen Urtexte (ebenso wie in der lateinischen Wittenberger und Jenaer Ausgabe) mehrere Sätze am Rande von oben nach unten, oder umgekehrt, gedruckt gewesen sind, die dann von den Ausgaben willkürlich bald hier bald da eingeschoben worden sind. Auch der Beschluß der Gebote ist verschieden gestellt. Die Weimarische Ausgabe hat in ihrer Einleitung eine Reconstruction der Urschrift angedeutet. Wir haben darnach den Stücken ihre Stellung am Schlusse der betreffenden Abschnitte angewiesen und sie dadurch kenntlich gemacht (ebenso wie die Weimarische Ausgabe), daß wir sie eingerückt haben drucken lassen. Diese Stücke müssen bei den ihnen vorhergehenden Abtheilungen (den Geboten selbst, ihrer Uebertretung und ihrer Erfüllung) von unten nach oben gedruckt gedacht werden.}

Anno 1518.

\section{Das erste Gebot}

\begin{schriftzitat}
Du sollst nicht andere Götter haben.
\end{schriftzitat}

Einen GOtt haben, das ist, einen haben, von dem er sich versieht, in allem Guten gefördert, in allem Bösen geholfen werden; das will der einige wahre GOtt selbst sein, und ist es auch.

\section{Das andere Gebot}

\begin{schriftzitat}
Du sollst den Namen Gottes, seines Herrn, nicht unnütz anziehen.
\end{schriftzitat}

Denn GOttes Namen sind auch diese: Wahrheit, Weisheit,\footnote{„Weisheit“ fehlt in der Weimarischen.} Gütigkeit, Kraft, und alles
\quellseite{684}
Guten Name ist sein Name; die soll ihm selbst niemand zuschreiben.

\section{Das dritte Gebot}

\begin{schriftzitat}
Du sollst den Feiertag heiligen.
\end{schriftzitat}

Die beste Feier ist, daß die Seele mit allem Thun desselbigen Tages deß gewarte, der da erfüllt die Hungrigen und Ledigen mit Gütern, Luc. 1, 53. Denn Feier, das heißt ledig sein.

\section{Das vierte Gebot}

\begin{schriftzitat}
Du sollst ehren deinen Vater und deine Mutter.
\end{schriftzitat}

Die Ehre steht nicht in Grüßen oder Neigen allein, sondern in Thun und Lassen alles, was ihr Wille oder Nothdurft ist.

\section{Das fünfte Gebot}

\begin{schriftzitat}
Du sollst nicht tödten.
\end{schriftzitat}

Nicht allein tödtet, der den Tod oder Schaden thut, sondern auch, der ihnen\footnote{Weimarische Ausgabe: „yn“; ebenso in dem gleich folgenden „kurzen Beschluß der zehn Gebote“: „yr yn“, das ist: ihr ihnen.} nicht wehrt und vorkommt, als Augustinus sagt.

\section{Das sechste Gebot}

\begin{schriftzitat}
Du sollst nicht ehebrechen.
\end{schriftzitat}

Die ärgste Frucht der Lust des Fleisches nennt er allein, zu beweisen, wie böse die ganze Lust sei, aus der solche Untugend kommt.

\section{Das siebente Gebot}

\begin{schriftzitat}
Du sollst nicht stehlen.
\end{schriftzitat}

Fremde Güter zu sich ziehen, ist eine Frucht des Geizes und Begierde; darum so wird hier der Baum mit den Früchten verboten, Geiz und was aus Geiz folgen mag.

\section{Das achte Gebot}

\begin{schriftzitat}
Du sollst nicht falsch Zeugniß reden wider deinen Nächsten.
\end{schriftzitat}

\section{Das neunte Gebot}

\begin{schriftzitat}
Du sollst nicht begehren deines Nächsten Weib, Tochter oder Magd.
\end{schriftzitat}

\section{Das zehnte Gebot}

\begin{schriftzitat}
Du sollst nicht begehren deines Nächsten Haus oder Gut.
\end{schriftzitat}

\subsection{Ein kurzer Beschluß der zehn Gebote}

Spricht Christus selber: „Was ihr wollt, daß euch die Menschen thun sollen, dasselbe thut ihr ihnen auch, das ist das ganze Gesetz und alle Propheten“, Matth. 7, 12. Denn niemand will Undank leiden für seine Wohlthat, oder seinen Namen einem andern lassen. Niemand will Hoffahrt gegen ihm erzeigt haben. Niemand will Ungehorsam, Zorn, Untreuschheit seines Weibes, Verraubung seiner Güter, Lügen, Trügen
\quellseite{685}
,\footnote{Weimarische: „Kriegen“. Diese Lesart scheint uns nach dem „rebellionem“ im Lateinischen zwar möglich, aber nicht wahrscheinlich, denn die Wittenberger und die Jenaer haben an der entsprechenden Stelle der „kurzen Form, die zehn Gebote, Glauben und Vater Unser zu betrachten“: „kriegen“.} Afterreden leiden; sondern Liebe und Freundschaft, Dank und Hülfe, Wahrheit und Treue erfinden von seinem Nächsten. Das gebieten aber alles die zehn Gebote.

Also lehren die Gebote.

\section{Uebertretung der zehn Gebote}

\subsection{Die Uebertretung des ersten Gebots}

Wer in seiner Widerwärtigkeit Zauberei, schwarze Kunst, Teufels Bundesgenossen sucht. Wer Briefe, Zeichen, Kräuter, Wörter, Segen und desgleichen gebraucht. Wer Wunschruthen, Schatzbeschwörungen, Kristallensehen, Mantelfahren, Milchstehlen übt. Wer sein Werk und Leben nach erwählten Tagen, Himmelszeichen, und der Weißsager Dünsten richtet. Wer sich selbst, sein Vieh,\footnote{Weim. Ausg.: „sein sich“ statt: sein Vieh.} Haus, Kinder, und allerlei Gut vor Wölfen, Eisen, Feuer, Wasser, Schaden, mit erdichteten Gebeten segnet und beschwört. Wer sein Unglück und Widerwärtigkeit dem Teufel oder bösen Menschen zuschreibt, und nicht mit Liebe und Lob alles Böse und Gute von GOtt alleine aufnimmt, und ihm wieder heimträgt mit Danksagen und völliger Gelassenheit. Wer GOtt versucht, und in unnöthige Fährlichkeit Leibes oder Seele sich gibt. Wer in seiner Frömmigkeit, Verstand oder andern geistlichen Gaben hoffärtig ist. Wer GOtt und die Heiligen, mit Vergessen der Seelenmoth, nur um zeitliches Nutzens willen ehrt.

\subsection{Die Uebertretung des andern Gebots}

Wer ohne Noth oder aus Gewohnheit leichtlich schwört. Wer falschen Eid schwört, oder auch sein Gelübde bricht. Wer Uebelthun gelobt oder schwört. Wer mit GOttes Namen flucht. Wer närrische Fabeln von Gott schwätzt, und die Worte der Schrift leichtfertig verkehrt. Wer GOttes Namen nicht anruft in seiner Widerwärtigkeit, und nicht gebenedeyet in Liebe und Leid, in Glück und Unglück. Wer Ruhm und Ehre und Namen sucht von seiner Frömmigkeit, Weisheit etc. Wer Gottes Namen anruft fälschlich, als die Ketzer und alle hoffärtigen Heiligen.

\subsection{Die Uebertretung des dritten Gebots}

Wer Fressen, Saufen, Spielen, Tanzen, Müßiggang, Unkeuschheit treibt. Wer Faulheit, Amt GOttes verschlafen, versäumen, Spazieren, unnütz Schwätzen übt. Wer ohne sonder[e] Noth arbeitet und handelt. Wer nicht betet, nicht Christi Leiden bedenkt, nicht seine Sünde bereuet, und Gnade begehrt, also nur mit Kleidern, Essen und Geberden äußerlich feiert.

\subsection{Die Uebertretung des vierten Gebots}

Wer sich Armuths, Gebrechens, Verachtung seiner Eltern schämt. Wer ihnen nicht ihre Nothdurft mit Speise und Kleidern versorgt; viel mehr wer ihnen flucht, schlägt, nachredet, haßt, und ungehorsam ist. Wer nicht von Herzen groß von ihnen hält, um GOttes Gebots willen. Wer die Gebote der christlichen Kirche nicht hält mit Fasten, Feyern etc. Wer Priester Stand unehrt, nachredet und beleidigt. Wer seine Herren und Obrigkeit nicht ehrt, treu und gehorsam ist, sie seien gut oder böse. Hierinne sind alle Ketzer, Abtrünnigen, Apostaten, Verbannten, Verstockten etc.

\subsection{Die Uebertretung des fünften Gebots}

Wer mit seinem Nächsten zürnt. Wer zu ihm sagt Racha, das sind allerlei Zorns und Hasses Zeichen. Wer zu ihm sagt fatue, du Narr, das sind allerlei Schandworte, Fluch, Lästerung, Nachreden, Richten, Urtheilen, Hohnspruch etc. Wer seines Nächsten Sünde oder Gebrechen rügt, und nicht bedeckt und entschuldiget. Wer seinen Feinden nicht vergibt, nicht für sie bittet, nicht freundlich ist, nicht wohlthut. Und hierinne sind alle Sünden des Zorns und Hasses, als Todtschläge, Krieg, Rauben, Brennen, Zanken, Hadern, Trauern des Nächsten Glücks, Freuen seines Unglücks. Wer nicht übt die Werke der Barmherzigkeit.

\subsection{Die Uebertretung des sechsten Gebots}

Wer Jungfrauen schwächt, ehebricht,\footnote{In der Weimarischen ist „ehebricht“ zweimal. Erlanger: Ehebruch.} Blutschanden, und dergleichen Unkeuschheit wirkt. Wer unnatürlicher Weise oder Personen (das
\quellseite{686}
sind stumme Sünden gebraucht. Wer mit sichtbaren Worten, Liedlein, Historien, Bildern, die böse Lust äugt\footnote{äugen = sehen lassen. Weimarische: eugent. Baldi und die Erlanger: übt.} oder zeigt. Wer mit Sehen, Greifen, vulgaren Gedanken sich reizt und befleckt. Wer die Ursache nicht meidet, als Freien, Tanzen, Müßigkeit, Faulheit, Schlafen, und Weibs- oder Mannspersonen Gemeinschaft. Wer mit übrigem Schmucke, Berden etc. andere zur Unkeuschheit reizt. Wer eines andern Keuschheit nicht hilft bewahren mit Rath und That.

\subsection{Die Uebertretung des siebenten Gebots}

Ueber die Dieberei und Räuberei ist auch der Wucher große Sünde. Wer falsches Gewicht und Maß braucht, oder böse Waare für gut ausgibt. Wer unrechte Erbgüter und Zins einnimmt. Wer verdienten Lohn vorbehält, und Schuld verleugnet. Wer seinem nächsten Thürigen nicht brot, oder leiht ohne allen Aufsatz. Alle, die geizig sind, und einen reich zu werden, und wie sonst fremd Gut behalten oder zu sich gebracht wird.

\subsection{Die Uebertretung des achten Gebots}

Wer vor Gericht die Wahrheit schweigt und unterdrückt. Wer schädlich leugt und betriegt. Item, alle schädlichen Schmeichler und Ehrenbläser, Zweizüngige, die Uneinigkeit machen und Unfrieden. Wer des Nächsten gut Leben, Werke und Worte übel auslegt und schmäht. Wer denselben bösen Zungen Statt gibt, hilft, und nicht widersteht.

\subsection{Die Uebertretung des neunten und zehnten Gebots}

Diese zwei letzten Gebote gehören nicht in die Beichte, sondern sind [zum] Ziel und Mal gesetzt, da wir hinkommen sollen, und täglich durch Buße dahin arbeiten, mit Hülfe und Gnade GOttes. Denn die böse Neigung stirbt nicht eher gründlich, daß Fleisch werde denn zu Pulver, und neu geschaffen.

Die fünf Sinne werden eingeschlossen im fünften und sechsten Gebote. Die sechs Werke der Barmherzigkeit im fünften und siebenten. Die sieben Todsünden, Hoffart etc. im ersten und andern. Unkeuschheit im sechsten. Zorn und Haß im fünften. Fraß im sechsten. Trägheit im dritten, und wohl in allen. Die frem-

den Sünden sind in allen Geboten; denn mit Heißen, Rathen und Hülfe, wider alle Gebote gesündigt kann werden. Die rufenden und stummen Sünden sind wider das fünfte, sechste und siebente Gebot etc.

\begin{verse}
Sein selbst eigene Liebe, und GOttes\\
und des Nächsten Verachtung.

Nimmt GOtt, das sein ist, und entzeucht ihm,\\
was ihr ist.

Nimmt dem Nächsten, das sein ist, und entbeut\\
ihm nicht, das ihr ist.

Also thut die Natur, ihr selbst gelassen,\\
durch Adams erste Sünde.
\end{verse}

\section{Die Erfüllung der zehn Gebote GOttes}

\subsection{Die Erfüllung des ersten Gebots}

GOttes Furcht und Liebe im rechten Glauben und fest vertrauen, ganz bloß, lauter, in allen Dingen gelassen stehen, sie seien böse oder gut.

\subsection{Die Erfüllung des andern Gebots}

Lob, Ehre, Gebenedeyung, und Anrufen GOttes Namen, und seinen eigenen Namen und Ehre ganz vernichten, daß allein Gott gepriesen sei, der allein alle Dinge ist und wirkt.

\subsection{Die Erfüllung des dritten Gebots}

Sich zu GOtt bereiten und Gnade suchen. Das geschieht mit Beten, Messe und Evangelii zuhören, und Christi Leiden bedenken und also geistlich zum Sacrament gehen. Denn dies Gebot fordert eine geistarme Seele, die da Ihres Nichtsein vor Gott opfert,\footnote{Dieser Satz lautet in der Weimarischen: „die da gar nicht sein vor got opffert“. In der Wittenberger und in der Jenaer ist dieser Satz so gedruckt in der kurzen Form: „die da jres (nicht sein) für Gott opffert“. Der Sinn ist: die das Ihre, nämlich „Nichtsein“, Gott opfert. Am Schlusse dieses Abschnitts ist es so ausgedrückt: „Die Liebe Gottes etc. entbeut ihm ihr bloßes Nichts.“} daß er GOtt sei, und in ihr sein Werk und Namen bekomme, nach den zwei ersten Geboten.

\subsection{Die Erfüllung des vierten Gebots}

Williger Gehorsam und Unterthänigkeit allerley Gewalt, um GOttes Wohlgefallen willen. Als der Apostel St. Petrus sagt [1. Ep. 2, 20.], ohne alles Widerbell(en), Klagen und Murmeln.

\quellseite{687}
\subsection{Die Erfüllung des fünften Gebots}

Geduld, Sanftmüthigkeit, Gütigkeit, Friedlichkeit, Barmherzigkeit, und allerdings ein süßes, freundliches Herz, ohne allen Haß, Zorn, Bitterkeit gegen einen jeglichen Menschen, auch den Feinden.

\subsection{Die Erfüllung des sechsten Gebots}

Keuschheit, Zucht, Schamhaftigkeit in Werken, Worten, Gebärden\footnote{Walch und die Erlanger: Begehren.} und Gedanken. Auch Mäßigkeit in Essen, Trinken, Schlafen, und alles, was der Keuschheit förderlich ist.

\subsection{Die Erfüllung des siebenten Gebots}

Armuth des Geistes, Mildigkeit, Willigkeit seine Güter zu leihen und geben, ohne allen Geiz und Begierde leben.

\subsection{Die Erfüllung des achten Gebots}

Eine friedsame, heilsame Zunge, die niemand schadet, und jedermann frommt, die die Uneinigen sühnet, die Verlästerten entschuldiget und versieht; das ist, Wahrheit und Einfältigkeit in Worten.

\subsection{Die neunte und zehnte Erfüllung}

Das ist vollkommene Keuschheit und Verachtung zeitlicher Lust und Güter gründlich, das allein in jenem Leben vollbracht wird, Amen.

\begin{verse}
Die Liebe GOttes und des Nächsten,\\
und sein selbst Verachtung.

Enthält sich der Güter und Namen GOttes,\\
und entbeut ihm ihr bloßes Nichts.

Enthält sich der Güter des Nächsten, und entbeut ihm das Ihre, und sich selbst.

Also thut die Gnade GOttes, durch\\
Christum unsern HErrn.
\end{verse}

Merke, es ist ein großer Irrthum, so jemand zum heiligen Sacrament geht, der sich verläßt auf das untreue Vortheil, daß er gebeichtet, gebetet, und ihm nicht bewußt ist einiger Todsünde. Denn die werden ihnen selbst das Gericht empfahen [1 Cor. 11, 29.], sintemal sie nicht rein noch würdig werden durch solche ihre Werke oder Gewissen; ja, durch solche vermessene Reinigkeit vielmehr unrein werden. Aber wer da glaubt und vertraut, er wolle Gnade und Reinigkeit am Sacrament erlangen, dieser Glaube und Vertrauen macht reine und würdige Empfaher, welcher sich nicht [ver]läßt auf die obgenannten Werke, sondern auf das reine, getreue und gütige Wort und Verheißen Christi, da er [Matth. 11, 28.] spricht: „Kommet zu mir alle, die ihr arbeitet und beschweret seid mit Sünden, ich will euch erquicken.“ In dieser tröstlichen Worte Vermessenheit und sicherem Vertrauen soll man zugeben, und die also zugeben, werden nicht zu Schanden. Denn wir aus unsern Werken nimmer gewiß werden mögen, ob wir ohne Todsünde seien, und muß doch gewiß ohne Todsünde sein, wer ohne Schaden zum Sacrament gehen will. Darum steht die Gewißheit auf dem Worte und Verheißen und Rufen des getreuen Seligmachers Christi, wie gesagt ist.

Das meint auch Paulus [1 Cor. 11, 28.], sprechend: „Ein jeglicher Mensch prüfe sich selbst, und esse dann von diesem Brod.“ Der prüft aber sich selbst recht, der eines andern Uebel vergißt, und nicht richtet, sondern sich selbst empfindet, daß er arbeitet, und beschweret ist mit viel Sünden und Gebrechen, und also begierig wird der Gnade und Hülfe Christi; das ist die alleredelste und nächste Bereitung zum Sacrament. Denn, als St. Augustinus sagt: Die Speise sucht nicht mehr, denn eine hungrige, ledige Seele, fleucht auch nicht, denn die da satt sind und hoffärtig, die sich unter einander richten und [ver]dammen, als die thaten, zu denen der Apostel diese Worte schreibt. Denn, sollte der Apostel uns in diesen Worten aufgelegt haben, daß wir erkunden sollten uns selbst, bis daß wir gewiß wären, ob wir ohne Todsünde wären, hätte er uns unmöglich Ding aufgelegt, und allesamt beraubt des heilsamen Sacraments. Darum ist genug, daß du nicht wissest eine Todsünde, gewisser, grober Weise, oder gewisses Vorsatzes zur Todsünde; das Hinterstellige befiehl in die Gnade GOttes, und laß deinem Glauben deine Reinigkeit sein, so bist du gewiß.
